\chapter{光滑代数簇上的微分算子层}

\section{代数簇上的微分算子}

设 $X$ 是特征零的代数闭域 $k$ 上的仿射代数簇。以 $\mathcal{O}_X$ 记 $X$ 的结构层，以 $R(X)$ 记其整体截面，即 $X$ 上正则函数环。则 $R(X)$ 是交换 $k$-代数，可如第一章 §5 那样定义 $R(X)$ 上 $k$-线性微分算子环 $D(X)$。称此环为 $X$ 上的\emph{微分算子环}。微分算子的阶数定义 $D(X)$ 上一个递增环滤过 $(D_p(X);\ p\in\mathbb{Z})$，它满足第一章 §3 开头的性质 (i)--(v)。

如第一章所讨论，在 $X=k^n$ 的情形已知 $D(X)=D(n)$ 是 $n$ 变量多项式系数微分算子环。

可把 $X$ 实现为某仿射空间 $k^n$（$n\in\mathbb{Z}_+$）中的闭子集。设 $I(X)$ 为 $k[X_1,X_2,\dots,X_n]$ 中所有在 $X$ 上消失的多项式构成的理想。则 $R(X)=k[X_1,X_2,\dots,X_n]/I(X)$，以 $r$ 记 $k[X_1,X_2,\dots,X_n]$ 到 $R(X)$ 上的限制同态。定义
$$A=\{T\in D(n)\mid T(I(X))\subset I(X)\}.$$
显然 $A$ 是 $D(n)$ 的子代数。

设 $T\in A$。则 $T$ 诱导 $k[X_1,X_2,\dots,X_n]/I(X)$ 的线性自同态 $\varphi(T)$。映射 $\varphi$ 是 $A$ 到 $R(X)$ 的所有线性自同态环的同态。显然 $A$ 配置微分算子阶数滤过是一个滤过环。

此外，$k[X_1,X_2,\dots,X_n]$ 是 $A$ 的子环。因此有如下交换图
$$\begin{tikzcd}
k[X_1,\dots,X_n] \arrow[r,"r"] \arrow[d,hook] & R(X) \arrow[d,hook]\\
A \arrow[r,"\varphi"'] & \End_k(R(X))
\end{tikzcd}$$

特别地，对任意多项式 $P\in k[X_1,\dots,X_n]$，$\varphi(P)$ 是 $X$ 上乘以 $r(P)$。设 $T\in A\cap D_p(n)$，$f_0,f_1,\dots,f_p$ 是 $R(X)$ 中元素的 $(p+1)$-元组。则可取 $P_0,P_1,\dots,P_p\in k[X_1,X_2,\dots,X_n]$ 使得 $r(P_i)=f_i$（$0\le i\le p$）。故
$$[[\cdots[[\varphi(T),f_0],f_1]\cdots,f_{p-1}],f_p]=\varphi([[\cdots[[T,P_0],P_1]\cdots,P_{p-1}],P_p])=0.$$
因此 $\varphi(T)$ 是 $X$ 上阶数 $\le p$ 的微分算子。由此可知 $\varphi:A\to D(X)$ 是与微分算子阶数滤过相容的环同态。

此外，
$$J(X)=\{T\in D(n)\mid T(k[X_1,X_2,\dots,X_n])\subset I(X)\}$$
是 $A$ 的双边理想。显然 $J(X)$ 含于 $\varphi$ 的核中。

\begin{lemma}
设 $T\in D(n)$。则下列条件等价：
\begin{enumerate}[label=(\roman*)]
  \item $T\in J(X)$；
  \item $T=\sum_I P_I\partial^I$，其中 $P_I\in I(X)$。
\end{enumerate}
\end{lemma}

\begin{proof}
若 $T$ 的系数 $P_I$ 在 $X$ 上消失，显然微分算子 $T$ 属于 $J(X)$。反之若 $T\in J(X)$，则 $P_0=T(1)$ 在 $X$ 上消失，即 $P_0\in I(X)$。假设 $|I|<m$ 时 $P_I\in I(X)$。则 $T'=\sum_{|I|<m}P_I\partial^I$ 属于 $J(X)$。因此 $T''=T-T'\in J(X)$。另一方面，对任意 $J\in\mathbb{Z}_+^n$，$|J|=m$，
$$T''(X^J)=\Bigl(\sum_{|I|\ge m}P_I\partial^I\Bigr)(X^J)=J!P_J$$
在 $X$ 上消失，即 $P_J\in I(X)$。故由 $m$ 上的归纳断定对所有 $I\in\mathbb{Z}_+^n$ 有 $P_I\in I(X)$。
\end{proof}

记 $D$ 为商环 $A/J(X)$。$A$ 按微分算子阶数的滤过诱导 $D$ 上的商环滤过 $(D_p;\ p\in\mathbb{Z})$，满足 I.3 开头的条件 (i)--(v)。

由于 $J(X)$ 在 $\varphi$ 的核中，它定义 $D$ 到 $D(X)$ 的同态 $\Phi$。显然 $\Phi$ 是与环滤过相容的 $D$ 到 $D(X)$ 的同态。

\begin{proposition}
态射 $\Phi:D\to D(X)$ 是滤过环的同构。
\end{proposition}

先证 $\Phi$ 是单射。设 $T\in A$ 使 $\varphi(T)=0$。这意味着 $\varphi(T)(P+I(X))=T(P)+I(X)=0$，即对任意 $P\in k[X_1,X_2,\dots,X_n]$ 有 $T(P)\in I(X)$。故 $T\in J(X)$，$\Phi$ 是单射。

为证满射，先作如下注记。

\begin{lemma}
设 $p\in\mathbb{Z}_+$，$P_I\in k[X_1,X_2,\dots,X_n]$（$I\in\mathbb{Z}_+^n$，$|I|\le p$）。则存在阶数 $\le p$ 的微分算子 $T\in D(n)$，使得对所有 $|I|\le p$ 的 $I\in\mathbb{Z}_+^n$ 有 $T(X^I)=P_I$。
\end{lemma}

\begin{proof}
显然 $p=0$ 时成立。设 $p>0$ 且结论对 $p-1$ 成立。由归纳假设存在阶数 $\le p-1$ 的微分算子 $T'$，使得对所有 $|I|\le p-1$ 有 $T'(X^I)=P_I$。对所有 $|I|=p$ 的 $I$，记 $T'(X^I)=Q_I\in k[X_1,X_2,\dots,X_n]$。显然 $\partial^J$（$|J|=p$）零化 $X^I$（$|I|\le p-1$），且当 $|I|=|J|=p$ 时 $\partial^J(X^I)=I!\delta_{I,J}$。这意味着若定义
$$T''=\sum_{|J|=p}\frac{P_J-Q_J}{J!}\partial^J,$$
则 $T''$ 零化 $X^I$（$|I|\le p-1$），且
$$T''(X^I)=\Bigl(\sum_{|J|=p}\frac{P_J-Q_J}{J!}\partial^J\Bigr)(X^I)=P_I-Q_I.$$
因此对 $|I|=p$ 有 $(T'+T'')(X^I)=P_I$，对 $|I|<p$ 有 $(T'+T'')(X^I)=T'(X^I)=P_I$。
\end{proof}

现断言：对任意 $T\in D_p(X)$ 和阶数 $\le p$ 的 $S\in D(n)$，若对 $|I|\le p$ 的 $I\in\mathbb{Z}_+^n$ 有 $T(r(X^I))=r(S(X^I))$，则 $T\circ r=r\circ S$。$p=0$ 时无需证明。设 $p>0$。则对 $1\le j\le n$，
\begin{align*}
[T,r(X_j)](r(X^I))&=T r(X_jX^I)-r(X_j)T(r(X^I))\\
&=r(S(X_jX^I))-r(X_jS(X^I))=r([S,X_j](X^I))
\end{align*}
对所有 $|I|\le p-1$ 成立，且 $[T,r(X_j)]$ 和 $[S,X_j]$ 的阶数 $\le p-1$。故由归纳假设 $[T,r(X_j)]\circ r=r\circ[S,X_j]$。特别地对任意 $I\in\mathbb{Z}_+^n$，
\begin{align*}
T(r(X_jX^I))&=[T,r(X_j)](r(X^I))+r(X_j)T(r(X^I))\\
&=r([S,X_j](X^I))+r(X_j)T(r(X^I))\\
&=r(S(X_jX^I))+r(X_j)\bigl[T(r(X^I))-r(S(X^I))\bigr].
\end{align*}
故由 $|I|$ 上的归纳得对所有 $I\in\mathbb{Z}_+^n$ 有 $T(r(X^I))=r(S(X^I))$，这证明断言。

设 $T$ 是 $X$ 上阶数 $\le p$ 的微分算子。则可选 $P_I\in k[X_1,X_2,\dots,X_n]$（$|I|\le p$）使得对所有 $|I|\le p$ 有 $T(r(X^I))=r(P_I)$。由上一条引理存在阶数 $\le p$ 的微分算子 $S\in D(n)$ 使得对所有 $|I|\le p$ 有 $S(X^I)=P_I$。这意味着对 $|I|\le p$ 有 $T(r(X^I))=r(S(X^I))$。

由前面的结果得 $T\circ r=r\circ S$。特别地 $r(S(I(X)))=T(r(I(X)))=0$，即 $S\in A$。显然 $\varphi(S)=T$，$\Phi$ 是满射。这完成 1.2 的证明。

\begin{corollary}
设 $X$ 是仿射代数簇。对任意 $p\in\mathbb{Z}_+$，$D_p(X)$ 作为 $R(X)$-模（对左乘和右乘）是有限生成的。
\end{corollary}

\begin{proof}
可设 $X$ 是 $k^n$ 中的闭子集。由 I.8.9 知 $X=k^n$ 时陈述成立。由于 $k[X_1,X_2,\dots,X_n]$ 是诺特环，$A\cap F_pD(n)$ 作为 $k[X_1,X_2,\dots,X_n]$-模（对左乘和右乘）有限生成，且对任意 $p\in\mathbb{Z}$，$D_p$ 作为 $R(X)$-模（对左乘和右乘）有限生成。结论由 1.2 得出。
\end{proof}

设 $f\in R(X)$，$f\neq 0$，$X_f=\{x\in X\mid f(x)\neq 0\}$ 是 $X$ 中相应的主开集。则 $X_f$ 是仿射簇，且 $R(X_f)=R(X)_f$。以 $r_f$ 记 $R(X)$ 到 $R(X_f)$ 的限制映射。

\begin{proposition}
设 $T\in D_p(X)$。则存在唯一的微分算子 $\bar T\in D_p(X_f)$ 使得如下交换图交换：
$$\begin{tikzcd}
R(X) \arrow[r,"T"] \arrow[d,"r_f"' ] & R(X) \arrow[d,"r_f"]\\
R(X_f) \arrow[r,"\bar T"'] & R(X_f)
\end{tikzcd}$$
\end{proposition}

先证 $\bar T$ 的唯一性。只需证明如下引理。

\begin{lemma}
设 $S\in D(X_f)$ 使得对任意 $g\in r_f(R(X))$ 有 $S(g)=0$。则 $S=0$。
\end{lemma}

\begin{proof}
对 $S$ 的阶数 $p$ 归纳证明。若 $p=0$，$S\in R(X_f)$，条件立即导致 $S=0$。现设 $p>0$。则 $S'=[S,f]\in D_{p-1}(X_f)$ 且它零化 $r_f(R(X))$。故由归纳假设 $S'=0$。这意味着 $S$ 与 $f$ 可交换。取 $h\in R(X_f)$。则存在 $n\in\mathbb{Z}_+$ 使得 $f^nh\in r_f(R(X))$。这意味着 $f^nS(h)=S(f^nh)=0$。由于 $f^{-1}\in R(X_f)$，断定 $S(h)=0$。因此 $S=0$。
\end{proof}

只需证明 $\bar T$ 的存在性。

先讨论 $X=k^n$ 的情形。由于 $D(n)$ 作为 $k$-代数由 $X_i,\partial_i$（$1\le i\le n$）生成，只需对 $T=\partial_i$（$1\le i\le n$）证明 $\bar T$ 的存在性。但导子 $\partial_i$ 扩充到有理函数域 $k(X_1,X_2,\dots,X_n)$ 且满足
$$\partial_i\Bigl(\frac{g}{f^m}\Bigr)=\frac{\partial_i(g)f-mg\,\partial_i(f)}{f^{m+1}}$$
对任意 $g\in k[X_1,X_2,\dots,X_n]$ 和 $m\in\mathbb{Z}_+$。因此它们诱导 $k[X_1,X_2,\dots,X_n]_f$ 的导子。这完成 $D(n)$ 存在性的证明。

只需证明一般情形下 $\bar T$ 的存在性。可设 $X$ 作为闭子集嵌入某 $k^n$。设 $P\in k[X_1,X_2,\dots,X_n]$ 是在 $X$ 上限制为 $f$ 的多项式，以 $U$ 记 $k^n$ 中 $P$ 的零点集的补集（仿射开集）。则 $X\cap U=X_f$。由 1.2 可找到 $S\in A\cap F_pD(n)$ 使得 $\varphi(S)=T$。此微分算子扩充为 $U$ 上阶数 $\le p$ 的微分算子 $\bar S$。

\begin{lemma}
设 $S\in A$。则 $\bar S$ 把 $I(X)_P$ 映入自身。
\end{lemma}

\begin{proof}
对 $S$ 的阶数 $p$ 归纳证明。$p=0$ 时显然。设 $p>0$。则 $S'=[S,P]\in A$ 且阶数 $\le p-1$。故由归纳假设 $\bar S'$ 把 $I(X)_P$ 映入自身。取 $Q\in I(X)$。则
$$\bar S\Bigl(\frac{Q}{P^m}\Bigr)=\bar S'\Bigl(\frac{Q}{P^{m+1}}\Bigr)+P\,\bar S\Bigl(\frac{Q}{P^{m+1}}\Bigr)$$
且由归纳假设
$$\bar S\Bigl(\frac{Q}{P^{m+1}}\Bigr)-P^{-1}\bar S\Bigl(\frac{Q}{P^{m}}\Bigr)\in I(X)_P$$
对任意 $m\in\mathbb{Z}_+$。由 $m$ 上的归纳得对任意 $m\in\mathbb{Z}_+$ 有 $\bar S\frac{Q}{P^m}\in I(X)_P$。
\end{proof}

因此 $\bar S$ 诱导 $k[X_1,X_2,\dots,X_n]_P/I(X)_P=R(X)_f=R(X_f)$ 的线性自同态。如 1.2 之前的讨论，这实际上是 $X_f$ 上的微分算子。此外在 $r_f(R(X))$ 上它与 $T$ 一致。因此构造了 $\bar T$。这完成 1.5 的证明。

设 $X$ 是仿射代数簇，$f\in R(X)$，$f\neq 0$。则由 1.5 有良定义的限制映射 $\rho_f:D(X)\to D(X_f)$。1.5 的唯一性部分蕴含 $\rho_f$ 是环同态，故有如下结果。

\begin{proposition}
映射 $\rho_f:D(X)\to D(X_f)$ 是滤过环的同态。
\end{proposition}

特别地，$\rho_f$ 是 $R(X)$-模（对左乘和右乘）的同态。

\begin{lemma}
设 $D(X)_f$ 是把 $D(X)$ 看作 $R(X)$-模（对左乘）的局部化。则态射 $\rho_f$ 诱导 $D(X)_f$ 到 $D(X_f)$ 上的同构 $\beta_f$。
\end{lemma}

\begin{proof}
先设 $X_f$ 在 $X$ 中稠密。此时自然映射 $r_f:R(X)\to R(X_f)$ 是单射。故 $\beta_f(\frac{1}{f^m}T)=\frac{1}{f^m}\bar T=0$（对某 $T\in D(X)$）蕴含对任意 $g\in R(X)$ 存在 $s\in\mathbb{Z}_+$ 使 $f^sT(g)=0$。因此对所有 $g\in R(X)$ 有 $T(g)=0$，即 $T=0$。由此可知 $\beta_f$ 是单射。

为证 $\beta_f$ 满射，只需证明对任意 $T\in D(X_f)$ 存在 $m\in\mathbb{Z}_+$ 使 $(f^mT)(R(X))\subset R(X)$。对 $T$ 的阶数 $p$ 归纳证明。$p=0$ 时显然。

设 $p>0$。记 $g_1,g_2,\dots,g_n$ 为 $k$-代数 $R(X)$ 的生成元。由归纳假设，存在 $m\in\mathbb{Z}_+$ 使得 $f^m[T,g_i](R(X))\subset R(X)$（$1\le i\le n$），且 $f^mT(1)\in R(X)$。这意味着若 $h\in R(X)$ 满足 $f^mT(h)\in R(X)$，则
$$f^mT(g_ih)=f^m[T,g_i](h)+f^m g_iT(h)\in R(X).$$
利用 $f^mT(1)\in R(X)$ 并对单项式 $g_1^{i_1}g_2^{i_2}\cdots g_n^{i_n}$ 的长度归纳，可见此关系蕴含 $f^mT(R(X))\subset R(X)$。故此时 $\beta_f$ 是同构。

现考虑一般情形。假设 $X_f$ 在 $X$ 中不稠密。断言此时存在 $f'\in R(X)$ 使得 $X_f$ 与 $X_{f'}$ 不交且它们的并在 $X$ 中稠密。首先，若 $X_f$ 在 $X$ 中不稠密，可找到 $a_1\in R(X)$ 使 $a_1\neq 0$ 且在 $X_f$ 上消失。这意味着 $X_f$ 与 $X_{a_1}$ 不交且 $X_{f+a_1}=X_f\cup X_{a_1}$。若 $X_{f+a_1}$ 在 $X$ 中不稠密可重复此构造，由于 $X$ 是诺特拓扑空间，经有限步构造序列 $a_1,a_2,\dots,a_s$，使得 $X_f,X_{a_1},X_{a_2},\dots,X_{a_s}$ 是 $X$ 中两两不交的主开集且并在 $X$ 中稠密。若令 $f'=a_1+a_2+\cdots+a_s$，它显然具有所需性质。

现断言 $D(X_{f+f'})=D(X_f)\oplus D(X_{f'})$。显然 $R(X_{f+f'})=R(X_f)\oplus R(X_{f'})$。设 $\chi,\chi'\in R(X_{f+f'})$ 分别为 $X_f$ 和 $X_{f'}$ 的特征函数。断言对任意 $T\in D_p(X_{f+f'})$ 有 $[T,\chi]=[T,\chi']=0$。$T$ 阶数 $\le 0$ 时成立。对 $T$ 的阶数归纳。设 $p>0$。若 $T$ 的阶数 $\le p$，由归纳假设断言对 $[T,\chi]$ 成立。因此
$$0=[[T,\chi],\chi']=[T,\chi]\chi'-\chi'[T,\chi]=-\chi T\chi'-\chi'T\chi$$
即 $\chi T\chi'=-\chi'T\chi$。右乘 $\chi$ 得 $\chi T\chi'=\chi'T\chi=0$。因此 $T\chi=(\chi+\chi')T\chi=\chi T\chi$，类似地 $\chi'T=\chi'T(\chi+\chi')=\chi'T\chi'$。由对称性也有 $\chi T=\chi T\chi$，最终导致 $[T,\chi]=0$。因此对 $g\in R(X_f)$ 有 $T(g)=T(\chi g)=\chi T(g)$，$T(R(X_f))\subset R(X_f)$。用 $\chi'$ 的类似论证得 $T(R(X_{f'}))\subset R(X_{f'})$。因此 $T$ 诱导 $X_f$ 和 $X_{f'}$ 上的微分算子 $S$ 和 $S'$，且 $T=S\oplus S'$。

由证明的第一部分，$\beta_{f+f'}:D(X)_{f+f'}\to D(X_f)\oplus D(X_{f'})$ 是同构。关于 $\chi$ 局部化得 $\beta_f:D(X)_f\to D(X_f)$ 是同构。
\end{proof}

设 $U$ 是 $X$ 中开集。以 $\mathcal{P}_U$ 记含于 $U$ 的所有主开集按包含排序的族。若 $V,W\in\mathcal{P}_U$ 且 $V\subset W$，有自然环同态 $r_{VW}:D(W)\to D(V)$。显然 $(D(V);r_{VW})$ 是环的逆系统。记其逆极限为 $D(U)$。显然 $\mathcal{D}_X:U\longmapsto D(U)$ 是 $X$ 上环的预层。由 1.9，这是左乘 $\mathcal{O}_X$-模的层。这蕴含如下结果。

\begin{proposition}
$\mathcal{D}_X$ 是 $X$ 上的环层。
\end{proposition}

称 $\mathcal{D}_X$ 为 $X$ 上\emph{局部微分算子层}。

\begin{theorem}
设 $X$ 是仿射簇，$\mathcal{D}_X$ 是 $X$ 上局部微分算子层。则对任意仿射开子集 $U\subset X$ 有 $\Gamma(U,\mathcal{D}_X)=D(U)$。
\end{theorem}

\begin{proof}
$U$ 是 $X$ 的主开集时显然。设 $U$ 是 $X$ 的任意仿射开子集。取 $f\in R(X)$ 使 $X_f\subset U$。则若记 $g=f|_U$，有 $U_g=X_f$。这意味着
$$\Gamma(U_g,\mathcal{D}_U)=D(U_g)=D(X_f)=\Gamma(X_f,\mathcal{D}_X).$$
此外这些同构与限制态射相容。由于主开集 $\{X_f\mid f\in R(X)\}$ 构成 $X$ 拓扑的基，含于 $U$ 的那些构成 $U$ 拓扑的基。此外，由于 $\mathcal{D}_X|_U$ 和 $\mathcal{D}_U$ 是 $U$ 上的层且在拓扑基上一致，它们相等。这意味着 $\Gamma(U,\mathcal{D}_X)=\Gamma(U,\mathcal{D}_U)=D(U)$。
\end{proof}

设 $X$ 是 $k$ 上任意代数簇。对 $X$ 中任意开集 $U$，以 $\mathcal{B}_U$ 记 $U$ 中所有仿射开子集按包含排序的族。若 $V,W\in\mathcal{B}_U$ 且 $V\subset W$，有自然环同态 $r_{VW}:D(W)\to D(V)$。显然 $(D(V);r_{VW})$ 是环的逆系统。记其逆极限为 $D(U)$。同样 $\mathcal{D}_X:U\longmapsto D(U)$ 是 $X$ 上环的预层。

\begin{proposition}
设 $X$ 是 $k$ 上的代数簇。则 $\mathcal{D}_X$ 是 $X$ 上的环层。
\end{proposition}

此结果以及 1.10 是如下引理的特例。设 $\mathcal{C}$ 是这样的范畴：其中任意对象的逆系统都有逆极限。设 $X$ 是拓扑空间，$\mathcal{B}$ 是 $X$ 拓扑的开集基。称 $\mathcal{B}$ 上取值于 $\mathcal{C}$ 的预层 $\mathcal{F}$ 为对象族 $\mathcal{F}(U)$（$U\in\mathcal{B}$）和态射族 $\rho_{VU}:\mathcal{F}(V)\to\mathcal{F}(U)$（对任意满足 $U\subset V$ 的对 $(U,V)$ 定义），满足
\begin{enumerate}[label=(\roman*)]
  \item 对任意 $U\in\mathcal{B}$，$\rho^U_U=\operatorname{id}$；
  \item 对任意 $U,V,W\in\mathcal{B}$ 且 $U\subset V\subset W$，$\rho^W_U=\rho^V_U\circ\rho^W_V$。
\end{enumerate}
则可通过令 $\mathcal{F}'(U)$ 等于所有满足 $V\subset U$ 的 $V\in\mathcal{B}$ 的 $\mathcal{F}(V)$ 的逆极限，定义 $X$ 上的预层 $\mathcal{F}'$。此外对任意 $U\in\mathcal{B}$ 有 $\mathcal{F}'(U)=\mathcal{F}(U)$。

\begin{lemma}
若 $\mathcal{F}$ 满足如下条件，则 $X$ 上的预层 $\mathcal{F}'$ 是 $X$ 上的层：

(F) 对 $U\in\mathcal{B}$ 的任意由 $U_i\in\mathcal{B}$ 构成的覆盖 $(U_i;\ i\in I)$ 以及任意对象 $T\in\mathcal{C}$，把 $f\in\operatorname{Hom}(T,\mathcal{F}(U))$ 映为族 $(\rho^{U_i}_U\circ f)_{i\in I}\in\prod_{i\in I}\operatorname{Hom}(T,\mathcal{F}(U_i))$ 的映射，是 $\operatorname{Hom}(T,\mathcal{F}(U))$ 到所有满足如下条件的族 $(f_i;\ i\in I)$ 的集合的双射：对任意指标对 $(i,j)$ 和任意 $V\in\mathcal{B}$ 且 $V\subset U_i\cap U_j$，有 $\rho^{V}_{U_j}\circ f_i=\rho^V_{U_i}\circ f_j$。
\end{lemma}

\begin{proof}
设 $\mathcal{B}'$ 是含于 $\mathcal{B}$ 的 $X$ 拓扑的另一个基。则可定义另一个预层 $\mathcal{F}''$，它由 $\mathcal{B}'$ 上 $\mathcal{F}$ 定义的预层附着而来。由逆极限的定义，对 $X$ 上任意开集 $U$ 存在 $\mathcal{F}'(U)$ 到 $\mathcal{F}''(U)$ 的典范态射。若 $U\in\mathcal{B}$，此典范态射是从 $\mathcal{F}(U)$ 到 $\mathcal{F}''(U)$ 的态射。断言此态射是同构。事实上由条件 (F)，$\mathcal{F}''(U)$ 到 $\mathcal{F}(V)$（$V\in\mathcal{B}'$，$V\subset U$）的典范态射通过 $\mathcal{F}(U)$ 分解。此外由泛性质，典范态射 $\mathcal{F}''(U)\to\mathcal{F}(U)$ 和前述态射 $\mathcal{F}(U)\to\mathcal{F}''(U)$ 两个顺序的复合都是恒等态射。这证明断言。另一方面，这也蕴含对 $X$ 中任意开集 $U$，态射 $\mathcal{F}''(U)\to\mathcal{F}''(V)=\mathcal{F}(V)$（$V\in\mathcal{B}$，$V\subset U$）满足逆系统 $(F(V);\ V\subset U,\ V\in\mathcal{B})$ 逆极限的条件，故 $\mathcal{F}''(U)=\mathcal{F}'(U)$。

现设 $U$ 是 $X$ 的开子集，$(U_i)$ 是 $U$ 的开覆盖，$\mathcal{B}'$ 是 $\mathcal{B}$ 中含于至少一个 $U_i$ 的元素构成的子族。显然 $\mathcal{B}'$ 是 $U$ 拓扑的基，故 $\mathcal{F}'(U)$（相应地 $\mathcal{F}'(U_i)$）是 $V\in\mathcal{B}'$ 且 $V\subset U$（相应地 $V\subset U_i$）的 $\mathcal{F}(V)$ 的逆极限。由逆极限的定义现在得出 $\mathcal{F}'$ 是层。
\end{proof}

设 $U$ 是 $X$ 中开集，$T\in\mathcal{D}_X(U)$。若对任意仿射开集 $V\subset U$，微分算子 $r^V_U(T)$ 阶数 $\le p$，则称 $T$ \emph{阶数 $\le p$}。这在 $\mathcal{D}_X(U)$ 上定义递增滤过 $F\mathcal{D}_X(U)$。

\begin{lemma}
$\mathcal{D}_X(U)$ 上的滤过 $F\mathcal{D}_X(U)$ 是穷尽的。
\end{lemma}

\begin{proof}
设 $T\in\mathcal{D}_X(U)$。由于 $U$ 拟紧，可找到由仿射开集构成的 $U$ 的有限开覆盖 $(U_i;\ 1\le i\le s)$。取 $p\in\mathbb{Z}$ 使 $T$ 在覆盖各元素上的限制阶数 $\le p$。设 $V$ 是 $U$ 的任意仿射开子集，$S=r^V_U(T)$。断言 $S$ 阶数 $\le p$。取 $f_0,f_1,\dots,f_p\in R(V)$。则 $R=[\cdots[[S,f_0],f_1],\dots,f_p]$ 是 $V$ 上的微分算子，其在 $V\cap U_i$（$1\le i\le s$）上的限制均为零。这意味着 $R=0$，$S$ 阶数 $\le p$。
\end{proof}

因此此滤过满足 I.3 开头的性质 (i)--(v)。显然这样得到 $X$ 上局部微分算子层 $\mathcal{D}_X$ 由 $k$-向量子层构成的滤过 $F\mathcal{D}_X$。称之为\emph{按微分算子阶数的滤过}。在 $X$ 上任意仿射开集 $U$ 上有 $F_p\mathcal{D}_X(U)=D_p(U)$（$p\in\mathbb{Z}$）。因此可考虑分次环层 $\Gr\mathcal{D}_X$。它是交换环层且 $\Gr_0\mathcal{D}_X=\mathcal{O}_X$。

\begin{theorem}
设 $X$ 是 $k$ 上的代数簇。则
\begin{enumerate}[label=(\roman*)]
  \item 层 $\mathcal{D}_X$ 是拟凝聚 $\mathcal{O}_X$-模（对左乘和右乘）；
  \item 层 $F_p\mathcal{D}_X$（$p\in\mathbb{Z}$）是凝聚 $\mathcal{O}_X$-模（对左乘和右乘）；
  \item 层 $\Gr_p\mathcal{D}_X$（$p\in\mathbb{Z}$）是凝聚 $\mathcal{O}_X$-模。
\end{enumerate}
\end{theorem}

\begin{proof}
由于断言是局部的，可设 $X$ 是仿射的。对左乘，(i) 由前述结果得出，(ii) 由 1.3 得出，(iii) 由 (ii) 得出。由于 $\Gr_p\mathcal{D}_X$ 上左乘和右乘定义相同的 $\mathcal{O}_X$-模结构，(iii) 成立。另一方面，
$$0\longrightarrow F_{p-1}\mathcal{D}_X\longrightarrow F_p\mathcal{D}_X\longrightarrow\Gr_p\mathcal{D}_X\longrightarrow 0$$
对右乘和左乘都是正合列，故由 $p$ 上的归纳得右乘的 (ii)。由于 $\mathcal{D}_X$ 是 $F_p\mathcal{D}_X$（$p\in\mathbb{Z}$）的直极限，(i) 成立。
\end{proof}

设 $X$ 是代数簇。对任意仿射开集 $U\subset X$ 记 $T_X(U)=\Der_k(R(U))$。由 I.8.2.(iii) 有 $D_1(U)=R(U)\oplus T_X(U)$。设 $V$ 是 $U$ 的仿射开子集。则对任意 $T\in T_X(U)$ 有 $r^V_U(T)(1)=T(1)=0$，故 $r^V_U(T)\in T_X(V)$，且限制映射与此直和分解相容。这意味着赋值 $U\longmapsto T_X(U)$ 定义 $X$ 中所有仿射开集构成的基 $\mathcal{B}$ 上的预层。记 $X$ 上相应的预层为 $\mathcal{T}_X$。由于 $F_1\mathcal{D}_X=\mathcal{O}_X\oplus\mathcal{T}_X$ 且 $F_1\mathcal{D}_X$ 和 $\mathcal{O}_X$ 都是层，$\mathcal{T}_X$ 是 $X$ 上的层。称之为 $X$ 的\emph{切层}。其在开集 $U\subset X$ 上的局部截面称为 $U$ 上的\emph{局部向量场}。显然它有自然的 $\mathcal{O}_X$-模结构，作为这样的模它同构于 $\Gr_1\mathcal{D}_X$。由前面的 (iii) 得出如下结果。

\begin{proposition}
设 $X$ 是 $k$ 上的代数簇。则
\begin{enumerate}[label=(\roman*)]
  \item $X$ 的切层 $\mathcal{T}_X$ 是凝聚 $\mathcal{O}_X$-模；
  \item $F_1\mathcal{D}_X=\mathcal{O}_X\oplus\mathcal{T}_X$。
\end{enumerate}
\end{proposition}

显然若 $T,T'$ 是 $U$ 上两个向量场，其换位子 $[T,T']$ 是 $U$ 上的向量场。因此 $\mathcal{T}_X$ 是 $k$ 上 Lie 代数的层。
\begin{annote}
相对微分算子环 $\mathcal D_{X/S}$ 的一阶部分正是 $\Omega_{X/S}=\Omega_{S/\mathbb Z}\otimes\mathcal O_X$ 的对偶：$\mathcal D_{X/S,\le1}=\mathcal O_X\oplus\Omega_{X/S}^{\vee}$ 只在光滑时成立（一般时左端是 $\mathcal O_X\oplus\mathcal T_{X/S}$ 的子层）。记 $\Omega_{S/R}$ 的定义与“模 $R$ 的微分”全称。[来源：Stacks Project Tag 07BK（Definition 10.131.2, module of Kähler differentials）；原书 §1, p55–62]
\end{annote}



\section{代数簇的光滑点}

设 $X$ 是特征零代数闭域 $k$ 上的代数簇。以 $\mathcal{O}_X$ 记其结构层。设 $x\in X$，以 $\mathcal{O}_x=\mathcal{O}_{X,x}$ 记 $\mathcal{O}_X$ 在 $x$ 处的茎。则 $\mathcal{O}_x$ 是诺特局部环，其极大理想 $\mathfrak{m}_x$ 由在 $x$ 处消失的函数芽构成。

\begin{lemma}
设 $x\in X$。则 $d(\mathcal{O}_{X,x})=\dim_x X$。
\end{lemma}

\begin{proof}
由于断言是局部的，可设簇 $X$ 是某 $k^n$ 的闭子集。则限制映射定义 $k[X_1,X_2,\dots,X_n]$ 到 $R(X)$ 上的满同态，其核 $I$ 由所有在 $X$ 上消失的多项式构成。可考虑 $k[X_1,X_2,\dots,X_n]$-模的正合列
$$0\longrightarrow I\longrightarrow k[X_1,X_2,\dots,X_n]\longrightarrow R(X)\longrightarrow 0$$
及其在 $x$ 处的局部化
$$0\longrightarrow I_x\longrightarrow k[X_1,X_2,\dots,X_n]_x\longrightarrow R(X)_x\longrightarrow 0.$$
这把 $\mathcal{O}_{k^n,x}=k[X_1,X_2,\dots,X_n]_x$ 模 $I_x$ 的商等同于 $\mathcal{O}_{X,x}=R(X)_x$。此外若以 $\mathcal{M}_x$ 记 $k[X_1,X_2,\dots,X_n]$ 中由 $X_i-x_i$（$1\le i\le n$）生成的极大理想，可见商态射把其局部化 $(\mathcal{M}_x)_x$ 映到 $\mathfrak{m}_x$ 上。这意味着 $\mathcal{O}_{X,x}$ 上的滤过 $(\mathfrak{m}_x^p;\ p\in\mathbb{Z}_+)$ 与 $\mathcal{O}_{X,x}$ 作为 $\mathcal{O}_{k^n,x}$-模的滤过 $((\mathcal{M}_x)^p_x\mathcal{O}_{X,x};\ p\in\mathbb{Z}_+)$ 一致。因此局部环 $\mathcal{O}_{X,x}$ 的维数 $d(\mathcal{O}_{X,x})$ 等于模 $\mathcal{O}_{X,x}$ 在局部环 $\mathcal{O}_{k^n,x}$ 上的维数。由 I.4.2 和 I.4.6 得
$$d(\mathcal{O}_{X,x})=\dim_x(\operatorname{supp}(R(X)))=\dim_x V(I)=\dim_x X.$$
\end{proof}

称向量空间 $T_x^*(X)=\mathfrak{m}_x/\mathfrak{m}_x^2$ 为 $X$ 在 $x$ 处的\emph{余切空间}，其线性对偶 $T_x(X)$ 为 $X$ 在 $x$ 处的\emph{切空间}。

\begin{proposition}
设 $x\in X$。则切空间 $T_x(X)$ 有限维且
$$\dim_k T_x(X)\ge\dim_x X.$$
\end{proposition}

\begin{proof}
这由前面的结果及 I.2.8 立即得出。
\end{proof}

设 $f\in\mathcal{O}_x$。则 $f-f(x)\in\mathfrak{m}_x$，以 $df(x)$ 记其在 $T_x^*(X)$ 中的像。

\begin{lemma}
线性映射 $d:\mathcal{O}_x\to T_x^*(X)$ 满足
$$d(fg)(x)=f(x)\,dg(x)+g(x)\,df(x)$$
对任意 $f,g\in\mathcal{O}_x$。
\end{lemma}

\begin{proof}
有
\begin{align*}
d(fg)(x)&=fg-f(x)g(x)+\mathfrak{m}_x^2\\
&=fg-f(x)g(x)-(f-f(x))(g-g(x))+\mathfrak{m}_x^2\\
&=g(x)(f-f(x))+f(x)(g-g(x))+\mathfrak{m}_x^2\\
&=f(x)\,dg(x)+g(x)\,df(x).
\end{align*}
\end{proof}

例如，若 $X=k^n$，则对任意芽 $f\in k[X_1,X_2,\dots,X_n]_x$ 有
$$df(x)=\sum_{i=1}^n(\partial_i f)(x)\,dX_i(x),$$
且 $(dX_1(x),dX_2(x),\dots,dX_n(x))$ 构成 $T_x^*(k^n)$ 的一组基。因此把向量 $(\xi_1,\xi_2,\dots,\xi_n)\in k^n$ 映为切向量 $f\longmapsto\sum_i\xi_i(\partial_i f)(x)$ 的映射是 $k^n$ 到 $T_x(k^n)$ 的同构。

设 $X$ 和 $Y$ 是 $k$ 上两个代数簇，$\varphi$ 是 $X$ 到 $Y$ 的态射。则对任意 $x\in X$，它诱导态射 $\varphi_x:\mathcal{O}_{Y,\varphi(x)}\to\mathcal{O}_{X,x}$，定义为 $\varphi_x(f)=f\circ\varphi$（$f\in\mathcal{O}_{Y,\varphi(x)}$）。显然 $\varphi_x(\mathfrak{m}_{\varphi(x)})\subset\mathfrak{m}_x$，这蕴含 $\varphi_x(\mathfrak{m}^2_{\varphi(x)})\subset\mathfrak{m}_x^2$，得到线性映射 $T^*_x(\varphi):T^*_{\varphi(x)}(Y)\to T^*_x(X)$。若 $f\in\mathcal{O}_{Y,\varphi(x)}$，有
$$T^*_x(\varphi)(df(\varphi(x)))=d(\varphi_x(f))(x)=d(f\circ\varphi)(x).$$
$T^*_x(\varphi)$ 的转置 $T_x(\varphi):T_x(X)\to T_{\varphi(x)}(Y)$ 称为 $\varphi$ 在 $x$ 处的\emph{切线性映射}。设 $\xi\in T_x(X)$，$f\in\mathcal{O}_{Y,\varphi(x)}$，则
$$\bigl(T_x(\varphi)(\xi)\bigr)(df(\varphi(x)))=\xi\bigl(T^*_x(\varphi)(df(\varphi(x)))\bigr)=\xi(d(f\circ\varphi)(x)).$$

\begin{lemma}\ 
\begin{enumerate}[label=(\roman*)]
  \item 设 $X,Y,Z$ 是代数簇，$\alpha:X\to Y$，$\beta:Y\to Z$ 是代数簇的态射。设 $x\in X$。则
  $$T_x(\beta\circ\alpha)=T_{\alpha(x)}(\beta)\circ T_x(\alpha).$$
  \item 设 $Y$ 是 $X$ 的子簇，$j:Y\to X$ 是典范包含。则对任意 $y\in Y$，$T_y(j):T_y(Y)\to T_y(X)$ 是单射。
\end{enumerate}
\end{lemma}

\begin{proof}
(i) 由定义得出。

(ii) 断言是局部的，故可设 $Y$ 在 $X$ 中闭且 $X$ 仿射。设 $I$ 是 $R(X)$ 中所有在 $Y$ 上消失的函数构成的理想。则 $\mathcal{O}_{Y,x}$ 是环 $R(X)/I$ 在 $x$ 处的局部化，由局部化的正合性，它是 $\mathcal{O}_{X,x}$ 的商。这意味着线性映射 $T^*_x(j):T^*_x(X)\to T^*_x(Y)$ 是满射，其转置 $T_x(j)$ 是单射。
\end{proof}

现设 $X$ 是某 $k^n$ 的闭子空间。由上述讨论和前一个引理 (ii)，自然包含 $j:X\to k^n$ 的切线性映射 $T_x(j)$ 把 $T_x(X)$ 等同于 $k^n$ 的线性子空间。以下结果精确刻画此子空间。如前以 $I$ 记 $k[X_1,X_2,\dots,X_n]$ 中所有在 $X$ 上消失的多项式的理想。

\begin{lemma}
对任意 $x\in X$ 有
$$T_x(X)=\Bigl\{(\xi_1,\xi_2,\dots,\xi_n)\in k^n\ \Big|\ \sum_{i=1}^n\xi_i(\partial_iP)(x)=0,\ P\in I\Bigr\}.$$
\end{lemma}

\begin{proof}
由定义和前一引理 (ii) 证明中的讨论，$T_x(X)$ 是 $T^*_x(j):k^n\to T^*_x(X)$ 的核的正交补。另一方面，$\ker T^*_x(j)=\{df(x)\mid f\in I_x\}$。$I_x$ 中任意芽 $f$ 是有理函数 $\frac{P}{Q}$ 的芽，其中 $Q(x)\neq 0$，$P\in I$。因此 $df(x)=\frac{1}{Q(x)}dP(x)$，且 $\{df(x)\mid f\in I_x\}=\{dP(x)\mid P\in I\}$。
\end{proof}

现考虑代数簇 $X$ 上的函数 $x\longmapsto\dim_kT_x(X)$。

\begin{proposition}
代数簇 $X$ 上的函数 $x\longmapsto\dim_k T_x(X)$ 是上半连续的。
\end{proposition}

\begin{proof}
断言是局部的，故可设 $X$ 是某 $k^n$ 的闭子空间。由前一结果可把切空间 $T_x(X)$ 等同于
$$\Bigl\{(\xi_1,\dots,\xi_n)\in k^n\ \Big|\ \sum_{i=1}^n\xi_i(\partial_iP)(x)=0,\ P\in I\Bigr\}.$$
若 $\dim T_x(X)=p$，则存在多项式 $P_1,P_2,\dots,P_{n-p}\in I$ 使得矩阵 $[(\partial_iP_j)(x)]$ 的秩为 $n-p$。这意味着在 $x$ 的某邻域 $U$ 中其秩等于 $n-p$。特别地对 $y\in U\cap X$ 有 $\dim_kT_y(X)\le p$。
\end{proof}

若 $\dim_kT_x(X)=\dim_xX$，则称点 $x\in X$ 是\emph{光滑的}。换句话说，$x\in X$ 光滑当且仅当局部环 $\mathcal{O}_x$ 是正则的。

\begin{theorem}
设 $X$ 是 $k$ 上的代数簇。则
\begin{enumerate}[label=(\roman*)]
  \item $X$ 的所有光滑点之集在 $X$ 中开且稠密；
  \item 光滑点 $x\in X$ 含于 $X$ 的唯一不可约分支中。
\end{enumerate}
\end{theorem}

\begin{proof}
第二陈述由 I.2.10 立即得出。记 $V_1,V_2,\dots,V_r$ 为 $X$ 的不可约分支，令 $Y=\bigcup_{i\neq j}V_i\cap V_j$。则 $Y$ 是 $X$ 的闭子集，其补在 $X$ 中稠密。此外由 (ii)，$Y$ 不含 $X$ 的光滑点。故可设 $X$ 是其不可约分支的不交并。这把证明归约到不可约簇的情形。

由前面的结果，光滑点之集在 $X$ 中开，证明归约为证明不可约仿射簇 $X$ 中光滑点的存在性。

设 $n=\dim X$。由 Noether 规范化引理，可找到 $f_1,f_2,\dots,f_n\in R(X)$ 使得由 $P\longmapsto P(f_1,\dots,f_n)$ 定义的 $k[X_1,X_2,\dots,X_n]$ 到 $R(X)$ 的同态是单射，且 $R(X)$ 在其像 $B$ 上是整的。几何上，这定义了 $X$ 到 $k^n$ 上的满有限态射 $p$。有理函数域 $R(X)$ 是 $B$ 的分式域 $L$（同构于 $k(X_1,X_2,\dots,X_n)$）的代数扩张。由于 $R(X)$ 是有限生成 $k$-代数，$R(X)$ 由 $R(X)$ 的有限多个元素在 $L$ 上生成。这意味着域 $R(X)$ 是 $L$ 的有限扩张。此外由于 $k$ 特征为零，由本原元定理，存在元素 $g\in R(X)$ 在 $L$ 上生成 $R(X)$。

先断言可设 $g\in R(X)$（原文本处疑有笔误）。设 $S=B-\{0\}$ 为乘法系，$S^{-1}R(X)$ 为相应的分式环。断言 $S^{-1}R(X)=R(X)$（此处指分式域）。取 $f\in S^{-1}R(X)$。由于 $R(X)$ 是 $L$ 的代数扩张，$f$ 在 $L$ 上代数，即存在 $b_1,b_2,\dots,b_n\in L$ 使
$$f^n+b_1f^{n-1}+\cdots+b_{n-1}f+b_n=0,\qquad b_n\neq 0.$$
这意味着
$$\frac{1}{f}=-\frac{1}{b_n}(f^{n-1}+b_1f^{n-2}+\cdots+b_{n-1})\in S^{-1}R(X).$$
因此 $S^{-1}R(X)$ 是含 $L$ 和 $R(X)$ 的域，故等于 $R(X)$（分式域）。由此可知任一本原元形如 $g=\frac{h}{b}$，其中 $h\in R(X)$，$b\in B-\{0\}$。这意味着 $h\in R(X)$ 也是本原元。

因此在以下固定本原元 $g\in R(X)$。设 $\Omega$ 为 $R(X)$ 的代数闭包，则 $\Omega$ 也是 $L$ 的代数闭包。设 $s$ 为 $R(X)$ 在 $L$ 上的次数。则 $g$ 在 $\Omega$ 的 $L$-自同构群 $\operatorname{Aut}_L(\Omega)$ 作用下的轨道由 $s$ 个元素 $g_0=g,g_1,\dots,g_{s-1}$ 组成。由于 $g$ 在 $B$ 上整且 $\operatorname{Aut}_L(\Omega)$ 固定 $B$，可见 $g_i$（$0\le i\le s-1$）在 $B$ 上整。令
$$V(\lambda_1,\lambda_2,\dots,\lambda_s)=\prod_{i<j}(\lambda_i-\lambda_j).$$
则对任意 $T\in\operatorname{Aut}_L(\Omega)$ 有
$$T\bigl(V(g_0,g_1,\dots,g_{s-1})^2\bigr)=V(g_0,g_1,\dots,g_{s-1})^2.$$
因此 $D=V(g_0,g_1,\dots,g_{s-1})^2$ 属于 $L$。此外它在 $B$ 上整。由于 $B$ 整闭，断定 $D\in B$。

显然 $1,g,g^2,\dots,g^{s-1}$ 构成向量空间 $R(X)$ 在 $L$ 上的一组基。因此对任意 $h\in R(X)$ 存在 $a_0,a_1,\dots,a_{s-1}\in L$ 使
$$h=\sum_{i=0}^{s-1}a_i g^i.$$
若令
$$h_j=\sum_{i=0}^{s-1}a_i g_j^i\quad(0\le j\le s-1),$$
可见 $h_0=h,h_1,\dots,h_{s-1}$ 在 $B$ 上整。由 Cramer 法则
$$a_j=\frac{\Delta_j}{V(g_0,g_1,\dots,g_{s-1})},$$
其中 $\Delta_j$ 是把 Vandermonde 行列式第 $j$ 列换为 $(h_0,h_1,\dots,h_{s-1})^T$ 所得行列式。这导出
$$a_j=\frac{\Delta_jV(g_0,g_1,\dots,g_{s-1})}{D}.$$
由于 $B$ 整闭，$a_jD\in L$，且 $\Delta_jV(g_0,g_1,\dots,g_{s-1})$ 在 $B$ 上整，断定 $a_jD\in B$（$0\le j\le s-1$）。这意味着
$$h=\sum_{i=0}^{s-1}a_iD\frac{g^i}{D}.$$
因此 $\frac{g^i}{D}$（$0\le i\le s-1$）生成 $R(X)$ 的一个含 $R(X)$ 的 $B$-子模（此处按原文）。

设 $Y=\{y\in k^n\mid D(y)\neq 0\}$，$X'=\{x\in X\mid (D\circ p)(x)\neq 0\}$。则 $p$ 把 $X'$ 映到 $Y$ 上。取 $y\in Y$ 和 $x\in X'$ 使 $p(x)=y$。由构造，$D$ 是 $g$ 在 $L$ 上的极小多项式 $\mu$ 的判别式。此外
$$\mu(X)=X^s+c_1X^{s-1}+\cdots+c_{s-1}X+c_s=\prod_{i=0}^{s-1}(X-g_i).$$
因此其系数 $c_j$（$1\le j\le s$）在 $B$ 上整，且由于 $B$ 整闭，它们属于 $B$。由于 $D$ 是 $g_i$ 的对称多项式，它是初等对称多项式 $c_1,c_2,\dots,c_s$ 的多项式。特别地 $D(y)$ 是多项式 $P(X)=X^s+c_1(y)X^{s-1}+\cdots+c_{s-1}(y)X+c_s(y)$ 的判别式。由于 $D(y)\neq 0$，此多项式在 $k$ 中有 $s$ 个不同的根。显然其在 $k$ 中的一个根是 $g(x)$。这意味着其导数 $P'$ 满足
$$P'(g(x))=sg(x)^{s-1}+(s-1)c_1(y)g(x)^{s-2}+\cdots+c_{s-1}(y)\neq 0.$$
因此若 $T\in T_x(X)$ 属于 $T_x(p)$ 的核，则
\begin{align*}
0&=T\bigl(d(g^s+(c_1\circ p)g^{s-1}+\cdots+(c_{s-1}\circ p)g+c_s\circ p)(x)\bigr)\\
&=P'(g(x))\,T(dg(x))+T\bigl(g(x)^{s-1}d(c_1\circ p)(x)+\cdots+g(x)d(c_{s-1}\circ p)(x)+d(c_s\circ p)(x)\bigr)\\
&=P'(g(x))\,T(dg(x))\\
&\quad +g(x)^{s-1}T_x(p)(T)(dc_1(y))+\cdots+g(x)T_x(p)(T)(dc_{s-1}(y))\\
&\quad +T_x(p)(T)(dc_s(y))\\
&=P'(g(x))\,T(dg(x)),
\end{align*}
这导致 $T(dg(x))=0$。现取任意 $h\in R(X)$。则由上述讨论，$h=\sum_{i=0}^{s-1}(b_i\circ p)g^i$，其中 $b_i\in R(Y)$。因此
\begin{align*}
T(dh(x))&=\sum_{i=0}^{s-1}g(x)^iT(d(b_i\circ p)(x))
=\sum_{i=0}^{s-1}g(x)^iT_x(p)(T)(db_i(y))=0.
\end{align*}
最后若 $f\in\mathcal{O}_{X,x}$，$f=\frac{h}{a}$，其中 $a,h\in R(X)$，$a(x)\neq 0$，则
$$df(x)=\frac{a(x)dh(x)-h(x)da(x)}{a(x)^2}.$$
因此对任意 $f\in\mathcal{O}_{X,x}$ 有 $T(df(x))=0$，从而断定 $T=0$，$T_x(p)$ 是单射。由于 $T_y(k^n)$ 是 $n$ 维的，可见 $\dim_kT_x(X)\le n$。由前面的结果得 $\dim_kT_x(X)=n$，$x$ 是光滑点。
\end{proof}

代数簇 $X$ 称为\emph{光滑的}，若其所有点都光滑。由前面的结果有如下命题。

\begin{proposition}
设 $X$ 是光滑代数簇。则其不可约分支等于其连通分支。
\end{proposition}

这特别蕴含函数 $x\longmapsto\dim_kX$ 在光滑簇 $X$ 上局部常值。

现在分析 $X$ 的光滑点邻域。断言如下结果。

\begin{theorem}
设 $X$ 是代数簇，$x\in X$ 是光滑点且 $\dim_xX=n$。则存在：
\begin{enumerate}[label=(\roman*)]
  \item $x$ 的仿射开邻域 $U$；
  \item $U$ 上的正则函数 $f_1,f_2,\dots,f_n$ 和向量场 $D_1,D_2,\dots,D_n$，使得 $D_i(f_j)=\delta_{ij}$（$1\le i,j\le n$）。
\end{enumerate}
\end{theorem}

\begin{proof}
由于断言是局部的，可设 $X$ 是嵌入某 $k^m$ 的光滑不可约仿射闭簇。设 $I$ 是 $A=k[X_1,X_2,\dots,X_m]$ 中所有在 $X$ 上消失的多项式的理想。由于 $\dim_kT_x(X)=\dim X=n$，可找到多项式 $P_{n+1},P_{n+2},\dots,P_m\in I$ 使得矩阵 $[(\partial_iP_j)(x)]$ 的秩为 $m-n$。这意味着此矩阵的秩在 $x\in k^m$ 的某邻域 $V$ 中等于 $m-n$，且
$$T_x(X)=\{(\xi_1,\dots,\xi_m)\in k^m\mid \textstyle\sum_{i=1}^m\xi_i(\partial_iP_j)(x)=0,\ n+1\le j\le m\}.$$

以 $J$ 记 $A$ 中由 $P_{n+1},P_{n+2},\dots,P_m$ 生成的理想。先断言 $J_x=I_x$。显然由定义 $J\subset I$。设 $Y$ 为 $J$ 在 $k^m$ 中所有零点的集合，则 $X\subset Y$。有
$$\dim_xY\ge\dim X=\dim_kT_x(X)=\dim_kT_x(Y)\ge\dim_xY.$$
这意味着 $\dim X=\dim_xY$。因此 $X$ 是 $Y$ 的一个不可约分支。另一方面，由于 $\dim_xY=\dim_kT_x(Y)$，$x$ 是 $Y$ 的光滑点且由前面的结果位于 $Y$ 的唯一不可约分支中。这意味着存在 $x$ 在 $k^m$ 中的邻域 $V'\subset V$ 不与 $Y$ 的任何其他不可约分支相交。因此断定 $J_x=I_x$。考虑局部环 $(A/J)_x$，其极大理想为 $\mathfrak{n}_x=\mathfrak{m}_x/J_x$，且对任意 $p\in\mathbb{Z}_+$ 有 $\mathfrak{n}^p_x=\mathfrak{m}^p_x/J_x$。因此局部环 $(A/J)_x$ 的维数等于其作为 $A_x$-模的维数。由 I.4.5 得 $d((A/J)_x)=\dim_x(V(J))=\dim_xY=\dim X=n$。另一方面有有限维向量空间的正合列
$$0\longrightarrow(J_x+\mathfrak{m}_x^2)/\mathfrak{m}_x^2\longrightarrow\mathfrak{m}_x/\mathfrak{m}_x^2\longrightarrow\mathfrak{n}_x/\mathfrak{n}_x^2\longrightarrow 0,$$
其中 $(J_x+\mathfrak{m}_x^2)/\mathfrak{m}_x^2=\{df(x)\mid f\in J_x\}$ 由 $dP_i(x)$（$n+1\le i\le m$）张成，且由上述等同 $\mathfrak{m}_x/\mathfrak{m}_x^2\cong k^m$。这意味着 $\dim_k(\mathfrak{n}_x/\mathfrak{n}_x^2)=n$，$(A/J)_x$ 是正则局部环。由 I.2.10 它是整环，故 $J_x$ 是素的。这最终导致 $J_x=I_x$。

这意味着 $A$-模 $I/J$ 的支集不含 $x$。特别地存在 $g\in A$ 使得由 $g$ 决定的 $k^m$ 中主开集 $V''\subset V'$ 与 $\operatorname{supp}(I/J)$ 不交。因此 $(I/J)_g=0$，$J_g=I_g$。

可找到多项式 $P_1,P_2,\dots,P_n\in A$ 使得矩阵 $[(\partial_iP_j)(x);\ 1\le i,j\le m]$ 可逆。因此必要时改变 $g$，还可设它在主开集 $V''$ 上可逆。以 $Q$ 记此矩阵的逆。则 $Q$ 的矩阵系数属于 $A_g$。因此在 $V''$ 上可定义微分算子 $\delta_i=\sum_{j=1}^mQ_{ij}\partial_j$（$1\le i\le n$）。显然它们满足
$$\delta_iP_j=\sum_{k=1}^mQ_{ik}\partial_kP_j=\delta_{ij}$$
对任意 $1\le j\le m$。由于任意 $f\in J_g$ 可表示为 $f=\sum_{j=n+1}^m h_jP_j$（$h_j\in A_g$），有
$$\delta_i(f)=\delta_i\Bigl(\sum_{j=n+1}^m h_jP_j\Bigr)=\sum_{j=n+1}^m(\delta_i(h_j)P_j+h_j\delta_i(P_j))=\sum_{j=n+1}^m\delta_i(h_j)P_j\in J_g,$$
即 $J_g=I_g$ 在 $\delta_i$（$1\le i\le n$）作用下不变。这意味着 $\delta_i$ 诱导 $U=X\cap V''$ 上的局部向量场，记为 $D_i$（$1\le i\le n$）。此外若以 $f_i$（$1\le i\le n$）记 $P_i$ 在 $U$ 上的限制，可见
$$D_i(f_j)=\delta_i(P_j)=\delta_{ij}.$$
\end{proof}

称 $(f_1,f_2,\dots,f_n;D_1,D_2,\dots,D_n)$ 为 $U\subset X$ 上的一个\emph{坐标系}。

\begin{lemma}
设 $X$ 是代数簇，$x\in X$ 是光滑点且 $\dim_xX=n$。则存在 $x$ 的仿射开邻域 $U$ 和 $U\subset X$ 上的坐标系 $(f_1,f_2,\dots,f_n;D_1,D_2,\dots,D_n)$，使得 $D_1,D_2,\dots,D_n$ 构成 $T_X(U)$ 作为自由 $\mathcal{O}_X(U)$-模的一组基。此外 $[D_i,D_j]=0$（$1\le i,j\le n$）。
\end{lemma}

\begin{proof}
由于任意光滑点位于 $X$ 的唯一不可约分支中，可设前一个定理中的邻域 $U$ 是不可约的。则 $\dim U=n$。以 $R(U)$ 记 $U$ 上有理函数域。由于 $U$ 是 $n$ 维的，$R(U)$ 在 $k$ 上的超越次数等于 $n$。设 $(f_1,f_2,\dots,f_n;D_1,D_2,\dots,D_n)$ 是 $U$ 上的坐标系。

断言 $f_1,f_2,\dots,f_n$ 在 $k$ 上代数无关。否则可找到非零且次数最小的多项式 $P\in k[X_1,X_2,\dots,X_n]$ 满足 $P(f_1,f_2,\dots,f_n)=0$。这意味着
$$0=D_i(P(f_1,f_2,\dots,f_n))=\sum_{j=1}^n(\partial_jP)(f_1,f_2,\dots,f_n)D_i(f_j)=(\partial_iP)(f_1,f_2,\dots,f_n),$$
由 $P$ 次数的最小性，$\partial_iP=0$（$1\le i\le n$）。由于 $k$ 特征为零，断定 $P$ 是常数多项式，显然不可能。

设 $K$ 是 $R(U)$ 中由 $f_1,f_2,\dots,f_n$ 生成的子域。则 $K$ 在 $k$ 上的超越次数也等于 $n$，$R(U)$ 是 $K$ 的代数扩张。

由于 $U$ 上向量场是 $R(U)$ 的导子，它扩充为 $R(U)$ 的导子。另一方面 $R(U)$ 是 $K$ 的代数扩张。故由隐函数微分法，此导子由其限制到 $K$ 唯一确定。可知 $U$ 上向量场完全由它限制到 $R(U)$ 中由 $f_1,f_2,\dots,f_n$ 生成的子环决定。

设 $T\in T_X(U)$，令 $g_i=T(f_i)$（$1\le i\le n$）。则
$$\Bigl(T-\sum_{i=1}^n g_iD_i\Bigr)(f_j)=T(f_j)-\sum_{i=1}^n g_iD_i(f_j)=0,$$
由上述讨论得 $T=\sum_{i=1}^ng_iD_i$，即 $D_1,D_2,\dots,D_n$ 生成 $T_X(U)$。另一方面若 $\sum_{i=1}^nh_iD_i=0$（$h_i\in R(U)$），则
$$0=\Bigl(\sum_{i=1}^nh_iD_i\Bigr)(f_j)=h_j\quad(1\le j\le n).$$
故 $\mathcal{O}_X(U)$-模 $T_X(U)$ 是自由的，$(D_1,D_2,\dots,D_n)$ 是其一组基。

最后对任意 $1\le i,j,k\le n$ 有
$$[D_i,D_j](f_k)=D_i(D_j(f_k))-D_j(D_i(f_k))=0,$$
这意味着 $[D_i,D_j]=0$。
\end{proof}

设 $x\in X$，$T\in T_{X,x}$。则 $T$ 决定局部环 $\mathcal{O}_{X,x}$ 的导子。显然对任意 $f\in\mathfrak{m}_x^2$ 有 $T(f)\in\mathfrak{m}_x$。此外对 $f\in\mathcal{O}_x$，有 $T(f)(x)=T(f-f(x))(x)$，结果只依赖于 $df(x)$。因此映射 $f\longmapsto T(f)(x)$ 通过 $T^*_x(X)$ 分解，定义切向量 $T(x)\in T_x(X)$，满足 $T(x)(df(x))=T(f)(x)$（$f\in\mathcal{O}_x$）。于是构造了 $T_{X,x}$ 到 $T_x(X)$ 的线性映射。显然它决定几何纤维 $\mathcal{O}_x/\mathfrak{m}_x\otimes_{\mathcal{O}_x}T_{X,x}$ 到 $T_x(X)$ 的线性映射。

\begin{proposition}
设 $x$ 是代数簇 $X$ 的光滑点。则 $\mathcal{O}_x/\mathfrak{m}_x\otimes_{\mathcal{O}_x}T_{X,x}$ 到 $T_x(X)$ 的典范映射是向量空间的同构。
\end{proposition}

\begin{proof}
设 $n=\dim_xX$。由前面的引理知 $T_{X,x}$ 是自由 $\mathcal{O}_x$-模。更确切地说，存在 $f_1,f_2,\dots,f_n\in\mathcal{O}_x$ 和 $D_1,D_2,\dots,D_n\in T_{X,x}$，满足 $D_i(f_j)=\delta_{ij}$（$1\le i,j\le n$），使得 $(D_1,D_2,\dots,D_n)$ 是自由 $\mathcal{O}_x$-模 $T_{X,x}$ 的一组基。这意味着 $D_1,D_2,\dots,D_n$ 在 $\mathcal{O}_x/\mathfrak{m}_x\otimes_{\mathcal{O}_x}T_{X,x}$ 中的像构成其作为 $k$-向量空间的基。

另一方面 $D_i(x)$ 满足 $D_i(x)(df_j(x))=\delta_{ij}$（$1\le i,j\le n$），故它们线性无关。由于切空间 $T_x(X)$ 是 $n$ 维的，断定 $(D_i(x),1\le i\le n)$ 是 $T_x(X)$ 的一组基，映射是双射。
\end{proof}

前面的两个结果蕴含如下定理。

\begin{theorem}
设 $X$ 是光滑代数簇。则切层 $\mathcal{T}_X$ 是有限秩局部自由 $\mathcal{O}_X$-模。对任意 $x\in X$，$\mathcal{T}_X$ 的几何纤维自然同构于 $T_x(X)$。
\end{theorem}

设 $X$ 是光滑代数簇，$T(X)=\{(x,\xi)\mid \xi\in T_x(X),\ x\in X\}$。想在 $T(X)$ 上定义自然的代数簇结构。

先设 $X$ 上切层 $\mathcal{T}_X$ 是自由 $\mathcal{O}_X$-模。设 $(T_1,T_2,\dots,T_n)$ 是 $\mathcal{T}_X$ 的一组基。则有 $X\times k^n$ 到 $T(X)$ 上的双射 $\varphi$，定义为
$$\varphi(x,\xi_1,\xi_2,\dots,\xi_n)=\Bigl(x,\sum_{i=1}^n\xi_iT_i(x)\Bigr).$$
通过要求 $\varphi:X\times k^n\to T(X)$ 是同构，可在 $T(X)$ 上定义代数簇结构。设 $(T'_1,T'_2,\dots,T'_n)$ 是自由 $\mathcal{O}_X$-模 $\mathcal{T}_X$ 的另一组基，$\varphi':X\times k^n\to T(X)$ 是相应的映射。则存在元素在 $R(X)$ 中的可逆矩阵 $Q$ 使得 $T'_i=\sum_{j=1}^nQ_{ji}T_j$，这意味着
$$\varphi'(x,\xi_1,\dots,\xi_n)=\Bigl(x,\sum_i\xi_iT'_i(x)\Bigr)=\Bigl(x,\sum_{i,j}\xi_iQ_{ji}(x)T_j(x)\Bigr)=\varphi\bigl(x,Q(x)(\xi_1,\dots,\xi_n)\bigr),$$
且 $\varphi'$ 是同构当且仅当 $\varphi$ 是同构。因此 $X$ 上的代数结构与 $\mathcal{T}_X$ 基的选择无关。

现考虑任意光滑代数簇 $X$。由前面的定理可找到 $X$ 的开覆盖 $(U_1,U_2,\dots,U_s)$ 使得 $\mathcal{T}_X|_{U_i}=\mathcal{T}_{U_i}$ 是自由 $\mathcal{O}_{U_i}$-模。显然 $T(X)$ 是 $T(U_i)$（$1\le i\le s$）的并，由上述讨论每个 $T(U_i)$ 有自然的代数簇结构。此外由于此结构与基的选择无关，可见在交 $T(U_i)\cap T(U_j)$ 上由 $T(U_i)$ 和 $T(U_j)$ 结构诱导的结构相同。这在 $T(X)$ 上定义了一个代数簇结构。称 $T(X)$ 为 $X$ 的\emph{切丛}。有自然映射 $i:X\to T(X)$ 和 $p:T(X)\to X$，定义为 $i(x)=(x,0)$ 和 $p(x,\xi)=x$（$\xi\in T_x(X)$，$x\in X$）。显然这些映射是代数簇的态射。此外有如下显然的结果。

\begin{proposition}
设 $X$ 是光滑代数簇。则
\begin{enumerate}[label=(\roman*)]
  \item 切丛 $T(X)$ 是光滑代数簇，且 $\dim_{(x,\xi)}T(X)=2\dim_xX$；
  \item 纤维化 $p:T(X)\to X$ 是局部平凡的。
\end{enumerate}
\end{proposition}

类似地，若 $X$ 是光滑簇，可在 $X$ 上定义 $\mathcal{O}_X$-模
$$\mathcal{T}^*_X=\mathcal{H}om_{\mathcal{O}_X}(\mathcal{T}_X,\mathcal{O}_X).$$
由于 $\mathcal{T}_X$ 是有限秩局部自由的，$\mathcal{T}^*_X$ 也是。这意味着其在 $x\in X$ 的几何纤维自然同构于余切空间 $T^*_x(X)$。设 $U\subset X$ 是开集，$f\in\mathcal{O}_X(U)$。则它定义 $\mathcal{H}om_{\mathcal{O}_X}(\mathcal{T}_X,\mathcal{O}_X)(U)=\operatorname{Hom}_{\mathcal{O}_X(U)}(\mathcal{T}_X(U),\mathcal{O}_X(U))$ 中由 $T\longmapsto T(f)$ 给出的元素，记为 $df$，称为 $f$ 的\emph{微分}。显然对任意 $x\in U$ 有
$$df(T)(x)=T(f)(x)=T(x)(df(x)),$$
故可把 $df(x)$ 看作由局部截面 $df$ 决定的 $\mathcal{T}^*_X$ 的几何纤维 $T^*_x(X)$ 中的元素。

若 $(f_1,f_2,\dots,f_n;D_1,D_2,\dots,D_n)$ 是充分小的仿射开集 $U$ 上的坐标系，由前面的结果 $(D_1,D_2,\dots,D_n)$ 构成 $\mathcal{T}_X(U)$ 作为自由 $\mathcal{O}_X(U)$-模的一组基。由此可知 $(df_1,df_2,\dots,df_n)$ 是自由 $\mathcal{O}_X(U)$-模 $\mathcal{T}^*_X(U)$ 的对偶基。

如切丛的情形，可定义 $T^*(X)=\{(x,\omega)\mid \omega\in T^*_x(X),\ x\in X\}$ 以及 $T^*(X)$ 上的代数簇结构，使得在带坐标系 $(f_1,f_2,\dots,f_n;D_1,D_2,\dots,D_n)$ 的充分小仿射开集 $U\subset X$ 上，$(df_1,df_2,\dots,df_n)$ 是自由 $\mathcal{O}_X(U)$-模 $\mathcal{T}^*_X(U)$ 的对偶基，且有同构 $U\times k^n\to T^*(U)\subset T^*(X)$，由
$$\varphi^*(x,\xi_1,\xi_2,\dots,\xi_n)=\Bigl(x,\sum_{i=1}^n\xi_i\,df_i(x)\Bigr)$$
给出。此簇称为 $X$ 的\emph{余切丛}。有自然映射 $\iota:X\to T^*(X)$ 和 $\pi:T^*(X)\to X$，定义为 $\iota(x)=(x,0)$ 和 $\pi(x,\omega)=x$（$\omega\in T^*_x(X)$，$x\in X$）。显然这些映射是代数簇的态射。

\begin{proposition}
设 $X$ 是光滑代数簇。则
\begin{enumerate}[label=(\roman*)]
  \item 余切丛 $T^*(X)$ 是光滑代数簇，且 $\dim_{(x,\omega)}T^*(X)=2\dim_xX$；
  \item 纤维化 $\pi:T^*(X)\to X$ 是局部平凡的。
\end{enumerate}
\end{proposition}

最后补充关于光滑点处正则函数芽的“Taylor 级数”的注记。

设 $X$ 是代数簇，$x\in X$ 是光滑点，其局部环 $\mathcal{O}_x$ 是正则局部环。若 $n=\dim_xX$，由前面的结果存在 $x$ 的仿射开邻域 $U$ 和 $U$ 上坐标系 $(f_1,f_2,\dots,f_n;D_1,D_2,\dots,D_n)$，使得 $f_1(x)=f_2(x)=\cdots=f_n(x)=0$ 且 $D_1(x),D_2(x),\dots,D_n(x)$ 构成 $T_x(X)$ 的一组基。因此 $df_1(x),df_2(x),\dots,df_n(x)$ 构成 $T^*_x(X)$ 的一组基，$f_1,f_2,\dots,f_n$ 在正则局部环 $\mathcal{O}_x$ 中定义一个坐标系。故有自然态射 $k[X_1,X_2,\dots,X_n]\to\mathcal{O}_x$，由 $P\longmapsto P(f_1,f_2,\dots,f_n)$ 给出。显然任何非常数项的多项式的像在 $\mathcal{O}_x$ 中可逆。因此这个态射扩充为局部环的态射 $\varphi:A=k[X_1,X_2,\dots,X_n]_0\to\mathcal{O}_x$。由于 $\mathcal{O}_x$ 是正则局部环，由 I.2.9 知 $\Gr\varphi$ 是 $\Gr A$ 到 $\Gr\mathcal{O}_x$ 的同构。由于 $A$ 的滤过是 Hausdorff 的，这意味着 $\varphi:A\to\mathcal{O}_x$ 是单射。故可把 $A$ 看作 $\mathcal{O}_x$ 的子环。自然同态 $A\to\mathcal{O}_x\to\mathcal{O}_x/\mathfrak{m}_x\cong k$ 是满射，因此其核 $A\cap\mathfrak{m}_x$ 是 $A$ 的极大理想，即等于 $\mathfrak{m}$。这意味着对任意 $p\in\mathbb{Z}_+$ 有 $\mathfrak{m}^p\subset A\cap\mathfrak{m}_x^p$。因此对所有 $p\in\mathbb{Z}_+$ 有自然映射 $A/\mathfrak{m}^p\to\mathcal{O}_x/\mathfrak{m}_x^p$，且图
$$\begin{tikzcd}
0 \arrow[r] & \Gr_pA \arrow[r] \arrow[d] & A/\mathfrak{m}^{p+1} \arrow[r] \arrow[d] & A/\mathfrak{m}^p \arrow[r] \arrow[d] & 0\\
0 \arrow[r] & \Gr_p\mathcal{O}_x \arrow[r] & \mathcal{O}_x/\mathfrak{m}_x^{p+1} \arrow[r] & \mathcal{O}_x/\mathfrak{m}_x^p \arrow[r] & 0
\end{tikzcd}$$
交换。由于两行正合且第一个竖直箭头是同构，若最后一个也是同构，由五引理中间箭头是同构。因此由 $p$ 上的归纳断定 $A/\mathfrak{m}^p\to\mathcal{O}_x/\mathfrak{m}_x^p$ 对所有 $p\in\mathbb{Z}_+$ 是同构。由此可知对所有 $p\in\mathbb{Z}_+$ 有 $\mathfrak{m}^p=A\cap\mathfrak{m}_x^p$。此外对任意 $p\in\mathbb{Z}_+$ 有 $A+\mathfrak{m}_x^p=\mathcal{O}_x$。故 $A$ 在 $\mathcal{O}_x$ 中稠密，$\mathfrak{m}_x$-进拓扑在 $A$ 上诱导 $\mathfrak{m}$-进拓扑。这意味着 $A$ 和 $\mathcal{O}_x$ 的完备化环同构。由于 $A$ 的完备化 $\widehat A$ 是形式幂级数环 $k[[X_1,X_2,\dots,X_n]]$，可见 $\mathcal{O}_x$ 可等同于 $k[[X_1,X_2,\dots,X_n]]$ 的子环。

设 $\widehat{\mathfrak{m}}$ 是局部环 $\widehat A=k[[X_1,X_2,\dots,X_n]]$ 的极大理想，即
$$\widehat{\mathfrak{m}}=\Bigl\{F\ \Big|\ F=\sum_{|I|\ge1}a_IX^I\Bigr\}.$$
这意味着
$$\widehat{\mathfrak{m}}^p=\Bigl\{F\ \Big|\ F=\sum_{|I|\ge p}a_IX^I\Bigr\}.$$
设 $\bar X_1,\bar X_2,\dots,\bar X_n$ 是 $X_1,X_2,\dots,X_n$ 在 $\widehat{\mathfrak{m}}/\widehat{\mathfrak{m}}^2$ 中的像。则立即看到 $(\bar X_i;1\le i\le n)$ 是 $\widehat A$ 中的坐标系。此外自然同态 $k[\bar X_1,\bar X_2,\dots,\bar X_n]$ 到 $\Gr\widehat A$ 是同构。因此 $\widehat A=k[[X_1,X_2,\dots,X_n]]$ 是正则局部环。

显然有 $\mathfrak{m}=A\cap\widehat{\mathfrak{m}}$ 和 $\mathfrak{m}^p\subset\widehat{\mathfrak{m}}^p$（$p\in\mathbb{Z}_+$）。$A$ 到 $\widehat A$ 的自然包含诱导 $\Gr A$ 到 $\Gr\widehat A$ 的同构。因此如前有交换图
$$\begin{tikzcd}
0 \arrow[r] & \Gr_pA \arrow[r] \arrow[d] & A/\mathfrak{m}^{p+1} \arrow[r] \arrow[d] & A/\mathfrak{m}^p \arrow[r] \arrow[d] & 0\\
0 \arrow[r] & \Gr_p\widehat A \arrow[r] & \widehat A/\widehat{\mathfrak{m}}^{p+1} \arrow[r] & \widehat A/\widehat{\mathfrak{m}}^p \arrow[r] & 0
\end{tikzcd}$$
且由归纳再次断定竖直箭头是同构。这意味着 $\mathfrak{m}^p=A\cap\widehat{\mathfrak{m}}^p$，即
$$\mathfrak{m}^p=\Bigl\{f\in A\ \Big|\ f=\sum_{|I|\ge p}a_IX^I\Bigr\}.$$

设 $T$ 是 $U$ 上向量场。则它诱导 $\mathcal{O}_x$ 的导子。由 $p$ 上的归纳可见 $T(\mathfrak{m}_x^p)\subset\mathfrak{m}^{p-1}_x$（$p\in\mathbb{N}$）。因此 $T$ 在 $\mathcal{O}_x$ 的 $\mathfrak{m}_x$-进拓扑中连续，并扩充为 $\mathcal{O}_x$ 完备化的连续导子。另一方面，多项式环 $k[X_1,X_2,\dots,X_n]$ 在形式幂级数环 $k[[X_1,X_2,\dots,X_n]]$ 中稠密，故 $k[[X_1,X_2,\dots,X_n]]$ 的任意连续导子完全由其在 $X_i$（$1\le i\le n$）上的作用决定。由于 $D_i(f_j)=\delta_{ij}$（$1\le i,j\le n$），这意味着在上述同构下 $D_i$ 对应于 $\partial_i$（$1\le i\le n$）。

任意形式幂级数 $F\in k[[X_1,X_2,\dots,X_n]]$ 可唯一写成其 Taylor 级数
$$F=\sum_{I\in\mathbb{Z}_+^n}\frac{(\partial^IF)(0)}{I!}\,X^I.$$
这与上述讨论一起立即得出如下结果。

\begin{lemma}
设 $x$ 是代数簇 $X$ 的光滑点，$(f_1,f_2,\dots,f_n;D_1,D_2,\dots,D_n)$ 是 $x$ 邻域中的坐标系。则对任意 $f\in\mathcal{O}_x$ 和 $p\in\mathbb{Z}_+$，下列条件等价：
\begin{enumerate}[label=(\roman*)]
  \item $f\in\mathfrak{m}_x^p$；
  \item $(D^If)(x)=0$ 对所有 $|I|<p$ 的 $I\in\mathbb{Z}_+^n$。
\end{enumerate}
特别地，$(D^If)(x)=0$ 对所有 $I\in\mathbb{Z}_+^n$ 蕴含 $f=0$。
\end{lemma}
\begin{annote}
光滑点 $x\in X$ 的定义是 $\Omega_{X/k}$ 在该处局部自由、且 $\dim T_xX=\dim_xX$；只有在此处才有 $\mathrm{gr}\mathcal D_X\cong\mathrm{Sym}\,\mathcal T_X$（讲义 2.x），从而特征簇落在 $T^*X$ 且维数比较可用。一般点上没有这个同构，只能退回 $\mathcal O_X$-凝聚语言，这正是第三章引入凝聚化函子的动机。[来源：原书 §2–§3, p62–76；译者按]
\end{annote}



\section{光滑簇上的微分算子层}

设 $X$ 是特征零代数闭域 $k$ 上的光滑代数簇。以 $\mathcal{D}_X$ 记 $X$ 上局部微分算子层，以 $F\mathcal{D}_X$ 记按微分算子阶数的滤过，以 $\Gr\mathcal{D}_X$ 记相应的分次环层。

先描述 $\Gr\mathcal{D}_X$ 的结构。设 $U$ 是 $X$ 中仿射开集。则由定义 $\Gamma(U,\mathcal{D}_X)=D(U)$。如 I.8，对任意 $p\in\mathbb{Z}_+$ 和 $T\in D_p(U)$ 可定义映射 $\sigma_p(T):R(U)^p\to R(U)$：
$$\sigma_p(T)(f_1,f_2,\dots,f_p)=[[\cdots[[T,f_1],f_2],\dots,f_{p-1}],f_p].$$
如 I.8.3 中所证，此映射是对称 $k$-多重线性的，且 $\sigma_p(T)=0$ 当且仅当 $T\in D_{p-1}(U)$。此外对任意 $1\le i\le p$，映射
$$f\longmapsto\sigma_p(T)(f_1,f_2,\dots,f_{i-1},f,f_{i+1},\dots,f_p)$$
是 $U$ 上的向量场。由于 $\sigma_p(T)$ 对称，为证此点可设 $i=p$。显然这是 $U$ 上阶数 $\le 1$ 的微分算子且在常数上为零，故由 1.16.(ii) 它是向量场。因此 $\sigma_p(T)(f_1,f_2,\dots,f_p)(x)$ 只依赖于 $f_i$ 在 $x$ 处的微分 $df_i(x)$（$1\le i\le p$）。由此可知可在 $U$ 的余切丛 $T^*(U)$ 上定义函数 $\Symb_p(T)$：
$$\Symb_p(T)(x,\omega)=\frac{1}{p!}\sigma_p(T)(f,f,\dots,f)(x),$$
其中 $f\in R(U)$ 满足 $df(x)=\omega$。

\begin{lemma}\ 
\begin{enumerate}[label=(\roman*)]
  \item 函数 $\Symb_p(T)$ 在 $T^*(U)$ 上正则；
  \item 对固定的 $x\in U$，函数 $\Symb_p(T)$ 是 $T^*_x(X)$ 上 $p$ 次齐次多项式。
\end{enumerate}
\end{lemma}

\begin{proof}
由于断言是局部的，由 2.10 可设 $U$ 充分小，使得 $U$ 上存在坐标系 $(f_1,f_2,\dots,f_n;D_1,D_2,\dots,D_n)$，且映射 $(x,\xi_1,\xi_2,\dots,\xi_n)\longmapsto(x,\sum_{i=1}^n\xi_i\,df_i(x))$ 是 $U\times k^n$ 到 $T^*(U)$ 上的同构。另一方面，
$$(x,\xi_1,\xi_2,\dots,\xi_n)\longmapsto\frac{1}{p!}\sigma_p(T)\Bigl(\sum_i\xi_if_i,\sum_i\xi_if_i,\dots,\sum_i\xi_if_i\Bigr)(x)$$
是 $U\times k^n$ 上的正则函数，这意味着 $\Symb_p(T)$ 在 $T^*(U)$ 上正则。第二陈述显然。
\end{proof}

称函数 $\Symb_p(T)$ 为微分算子 $T$ 的 \emph{$p$-主符号}。

设 $\pi:T^*(X)\to X$ 是自然投射，定义为 $\pi(x,\omega)=x$（$\omega\in T^*_x(X)$，$x\in X$）。由于 $\pi$ 是局部平凡纤维化且 $x\in X$ 处的纤维是 $T^*_x(X)$，可见 $T^*_x(X)$ 上多项式的齐次次数的自然分次在直像层 $\pi_*(\mathcal{O}_{T^*(X)})$ 上诱导分次环层结构。显然符号映射 $\Symb_p$ 定义层 $F_p\mathcal{D}_X$ 到 $\Gr_p\pi_*(\mathcal{O}_{T^*(X)})$ 的同态，同名记之。它在子层 $F_{p-1}\mathcal{D}_X$ 上为零，故它定义层 $\Gr_p\mathcal{D}_X$ 到 $\pi_*(\mathcal{O}_{T^*(X)})$ 的第 $p$ 次齐次分量的同态。设 $\Symb:\Gr\mathcal{D}_X\to\pi_*(\mathcal{O}_{T^*(X)})$ 为相应的分次层同态。

\begin{theorem}
符号映射 $\Symb:\Gr\mathcal{D}_X\to\pi_*(\mathcal{O}_{T^*(X)})$ 是分次 $\mathcal{O}_X$-代数层的同构。
\end{theorem}

此结果的证明分几步。先证明符号映射是 $k$-代数层的同态。

\begin{lemma}
设 $U$ 是 $X$ 的开子集，$T,S\in\mathcal{D}_X(U)$ 分别是阶数 $\le p$ 和 $\le q$ 的算子。则
$$\Symb_{p+q}(TS)=\Symb_p(T)\,\Symb_q(S).$$
\end{lemma}

\begin{proof}
设 $f\in\mathcal{O}_X(U)$，定义映射 $\tau:\mathcal{D}_X(U)\to\mathcal{D}_X(U)$ 为 $\tau(T)=[T,f]$。则
$$\tau(TS)=[TS,f]=TSf-fTS=[T,f]S+T[S,f]=\tau(T)S+T\tau(S).$$
故对任意 $k\in\mathbb{Z}_+$，
$$\tau^k(TS)=\sum_{i=0}^k\binom{k}{i}\tau^{k-i}(T)\,\tau^i(S).$$
这意味着若固定 $x\in X$ 和 $\omega\in T^*_x(X)$ 使 $df(x)=\omega$，有
\begin{align*}
\Symb_{p+q}(TS)(x,\omega)&=\frac{1}{(p+q)!}\sigma_{p+q}(T)(f,f,\dots,f)(x)=\frac{1}{(p+q)!}\tau^{p+q}(TS)(x)\\
&=\frac{1}{p!\,q!}\tau^p(T)(x)\,\tau^q(S)(x)=\Symb_p(T)(x,\omega)\,\Symb_q(S)(x,\omega).
\end{align*}
\end{proof}

显然，由于纤维化 $\pi:T^*(X)\to X$ 是局部平凡的，分次环层 $\pi_*(\mathcal{O}_{T^*(X)})$ 的第零次齐次分量等于 $\mathcal{O}_X$，且符号映射 $\Symb_0$ 是恒等映射。另一方面，第一个齐次分量自然同构于 $\mathcal{T}_X$。此外对 $x$ 邻域上的任意向量场 $T$ 和正则函数 $f$ 有
$$\Symb_1(T)(x,df(x))=[T,f](x)=T(f)(x)=T(x)(df(x)).$$
由于由 1.16.(ii) 有 $F_1\mathcal{D}_X=\mathcal{O}_X\oplus\mathcal{T}_X$，断定 $\Symb_1$ 是 $\Gr_1\mathcal{D}_X$ 到 $\pi_*(\mathcal{O}_{T^*(X)})$ 第一个齐次分量的同构。由 $\pi$ 的局部平凡性，分次环层 $\pi_*(\mathcal{O}_{T^*(X)})$ 由其第零和第一次齐次分量生成。因此符号映射的像等于 $\pi_*(\mathcal{O}_{T^*(X)})$。只需证明其核为零。

\begin{lemma}
设 $T\in F_p\mathcal{D}_X(U)$。则 $\Symb_p(T)=0$ 当且仅当 $T$ 阶数 $\le p-1$。
\end{lemma}

\begin{proof}
断言是局部的，故可设 $U$ 仿射。对 $p$ 归纳证明。$p=0$ 时显然，故设 $p>0$。固定 $f\in\mathcal{O}_X(U)$，则 $[T,f]$ 是阶数 $\le p-1$ 的微分算子。取 $x\in U$ 和 $\omega\in T^*_x(X)$，令 $\eta=df(x)$。必要时缩小 $U$，可取 $g\in\mathcal{O}_X(U)$ 使 $dg(x)=\omega$。对任意 $h\in\mathcal{O}_X(U)$ 定义映射 $\tau_h:\mathcal{D}_X(U)\to\mathcal{D}_X(U)$ 为 $\tau_h(T)=[T,h]$。则对任意 $\lambda\in k$，
$$\tau_{f+\lambda g}(T)=[T,f+\lambda g]=[T,f]+\lambda[T,g]=\tau_f(T)+\lambda\tau_g(T).$$
由于 $\tau_f$ 和 $\tau_g$ 可交换，对任意 $k\in\mathbb{Z}_+$ 有
$$\tau_{f+\lambda g}^k(T)=\sum_{i=0}^k\binom{k}{i}\lambda^i\tau_g^i(\tau_f^{k-i}(T)).$$
由假设，映射
$$\lambda\longmapsto\Symb_p(T)(x,\eta+\lambda\omega)=\frac{1}{p!}\tau_{f+\lambda g}^p(T)(x)$$
在 $k$ 上恒为零。由于 $k$ 无限，对 $1\le i\le p$ 有 $\tau_g^{p-i}(\tau_f^i(T))(x)=0$。特别地，
$$\Symb_{p-1}([T,f])(x,\omega)=\frac{1}{(p-1)!}\tau_g^{p-1}(\tau_f(T))(x)=\frac{1}{(p-1)!}\tau_g^{p-1}([T,f])(x)=0$$
对任意 $\omega\in T^*_x(X)$ 成立。由于 $x\in U$ 任意，由归纳假设 $[T,f]$ 阶数 $\le p-2$。这意味着 $T$ 阶数 $\le p-1$。
\end{proof}

这完成定理 3.2 的证明。

\begin{proposition}
光滑簇 $X$ 上局部微分算子层 $\mathcal{D}_X$ 是局部自由 $\mathcal{O}_X$-模（对左乘和右乘）。

更确切地说，每点 $x\in X$ 有仿射开邻域 $U$ 和 $U$ 上的坐标系 $(f_1,f_2,\dots,f_n;D_1,D_2,\dots,D_n)$，使得
\begin{enumerate}[label=(\roman*)]
  \item 对任意 $I,J\in\mathbb{Z}_+^n$ 有 $D^I\circ D^J=D^{I+J}$；
  \item $(D^I;\ I\in\mathbb{Z}_+^n,\ |I|\le p)$ 是自由 $\mathcal{O}_X(U)$-模 $F_p\mathcal{D}_X(U)$（对左乘和右乘）的一组基；
  \item $(D^I;\ I\in\mathbb{Z}_+^n)$ 是自由 $\mathcal{O}_X(U)$-模 $\mathcal{D}_X(U)$（对左乘和右乘）的一组基。
\end{enumerate}
\end{proposition}

\begin{proof}
设 $U$ 是 $x$ 的邻域，$(f_1,f_2,\dots,f_n;D_1,D_2,\dots,D_n)$ 是如 2.10 中 $U$ 上的坐标系。则对 $1\le i,j\le n$ 有 $[D_i,D_j]=0$，(i) 成立。

记 $\xi_i=\Symb_1(D_i)$（$1\le i\le n$）。则 $\pi_*(\mathcal{O}_{T^*(X)})(U)$ 是以 $(\xi^I;\ I\in\mathbb{Z}_+^n)$ 为基的自由 $\mathcal{O}_X(U)$-模，其齐次分量是自由 $\mathcal{O}_X(U)$-模。由正合列
$$0\longrightarrow F_{p-1}\mathcal{D}_X\longrightarrow F_p\mathcal{D}_X\longrightarrow\Gr_p\mathcal{D}_X\longrightarrow 0$$
及定理 3.2，由 $p$ 上的归纳断定 $F_p\mathcal{D}_X(U)$ 是自由 $\mathcal{O}_X(U)$-模，且由 $(D^I;\ I\in\mathbb{Z}_+^n,\ |I|\le p)$ 生成。若 $\sum_{|I|\le p}f_ID^I=0$，取 $p$ 次主符号得 $|I|=p$ 时 $f_I=0$ 且 $\sum_{|I|\le p-1}f_ID^I=0$。故由递降归纳得 $|I|\le p$ 时 $f_I=0$。这证明 (ii) 和 (iii)。
\end{proof}

\begin{proposition}
设 $X$ 是特征零代数闭域上的光滑仿射簇。则
\begin{enumerate}[label=(\roman*)]
  \item $\Gr D(X)$ 是诺特环；
  \item $\Gr D(X)$ 是由 $\Gr_1D(X)$ 生成的 $R(X)$-代数。
\end{enumerate}
\end{proposition}

\begin{proof}
由于 $\pi:T^*(X)\to X$ 是局部平凡纤维化且纤维是向量空间，断定 $\pi$ 是仿射态射。这意味着 $T^*(X)$ 是仿射簇。故
$$\Gr D(X)=\Gamma(X,\Gr\mathcal{D}_X)\cong\Gamma(X,\pi_*(\mathcal{O}_{T^*(X)}))=\Gamma(T^*(X),\mathcal{O}_{T^*(X)})=R(T^*(X))$$
是有限生成 $k$-代数和诺特环。此外由于 $\pi_*(\mathcal{O}_{T^*(X)})$ 作为 $\mathcal{O}_X$-代数由第一次齐次分量生成，对任意 $p\in\mathbb{Z}_+$，$\pi_*(\mathcal{O}_{T^*(X)})_1\otimes_{\mathcal{O}_X}\pi_*(\mathcal{O}_{T^*(X)})_p\to\pi_*(\mathcal{O}_{T^*(X)})_{p+1}$ 是满射。由于 $X$ 仿射，这意味着整体截面的相应态射是满射。因此 $\Gr D(X)$ 作为 $R(X)$-代数由 $\Gr_1D(X)$ 生成。
\end{proof}

\begin{theorem}
设 $X$ 是特征零代数闭域上的光滑仿射簇。则
\begin{enumerate}[label=(\roman*)]
  \item $D(X)$ 是左诺特且右诺特的环；
  \item 环 $D(X)$ 由 $R(X)$ 和 $X$ 上整体向量场生成。
\end{enumerate}
\end{theorem}

\begin{proof}
由前面的命题，$D(X)$ 是满足 I.3 开头条件 (i)--(vii) 的滤过环。故 (i) 由 I.3.4 得出。

设 $A$ 是 $D(X)$ 中由 $R(X)$ 和 $X$ 上整体向量场生成的子环。设 $F A$ 是 $A$ 上的诱导滤过。则有 $\Gr A$ 到 $\Gr D(X)$ 的单射同态，由前面的命题 (ii) 它也是满射。这意味着 $A=D(X)$。
\end{proof}
